\section{The anomaly catalog}
\label{sec:anomalies}

Four anomaly predicates over histories are formalized in \TLA{}
(\texttt{Anomalies.tla}); each has a TLC witness at $|A|{=}2$,
$|\fld{Cells}|\le 2$, $\fld{MaxOps}\le 4$, reproduced by the artifact's
\texttt{reproduce.sh} and reproduced independently for this version.
\Cref{tab:classical} names each predicate's classical counterpart and the
regime feature that forces the port. The numbering ($A_1, A_2, A_3, A_6$) is
kept from the public artifact; an earlier $A_5$ was subsumed by \A1.
\Cref{sec:residue} then defines, outside the catalog, the predicate for what
the four leave behind.

\begin{table}[t]
\centering
\caption{Each cataloged predicate, its classical original, and what the agent
regime changes. \A4 (split view) needs replication and is outside the
single-store model.}
\label{tab:classical}
\footnotesize
\setlength{\tabcolsep}{3pt}
\begin{tabular}{@{}p{1.55cm}p{2.75cm}p{3.85cm}@{}}
\toprule
Predicate & Classical condition & What the port changes \\
\midrule
\A1 stale generation & failed backward validation~\cite{KungRobinson1981OCC}; its absence is commit-order read consistency (\cref{thm:corc}) & the read phase is an inference call: unbounded, unlocked, not cheaply re-validated \\
\A2 phantom tool & dangling capability reference~\cite{Miller2006Robust} & the registry is mutable state with no relational analogue; the call reaches whatever is live at dispatch \\
\A3 causal cascade & recoverability, externalization in the role of commit~\cite{BernsteinHadzilacosGoodman1987} & effects are irreversible, so the dependent must be \emph{held}, not aborted after the fact \\
\A6 effect reordering & order preservation of external effects (2PC, write buffering) & a batch the model declared parallelizable is dispatched concurrently; completion order is the commit order \\
\Rp{} (residue, \cref{sec:residue}; not an anomaly) & view maintenance~\cite{GuptaMumick1995ViewMaint}, with the view outside the system; why-provenance~\cite{BunemanEtAl2001Why} for the basis & the effect is irreversible and its basis is revisable, so the window after commit never closes; only an obligation can be guaranteed (\cref{prop:owed}) \\
\bottomrule
\end{tabular}
\end{table}

\subsection{\A1: stale generation}
\label{sec:a1}
\begin{definition}[Stale generation]
\label{def:a1}
A history $h$ exhibits \emph{stale generation} when there exist
$i\neq j$ with $h[i].\fld{agent}\neq h[j].\fld{agent}$ and
$c\in h[i].\rs\cap h[j].\ws$ such that
\[
  \begin{gathered}
  h[i].\rt < h[j].\wt < h[i].\wt\\
  \text{and}\quad h[i].\rv[c] \neq h[j].\wv[c].
  \end{gathered}
\]
\end{definition}

\noindent\emph{Witness.} TLC finds a five-state counterexample: $a_2$ reads
$c_1=\mathrm{NULL}$ at $\rt=0$; $a_1$ commits $c_1=v_1$ at $\wt=1$; $a_2$
commits at $\wt=2$ with $\rv[c_1]=\mathrm{NULL}$.

\smallskip\noindent\emph{Classical placement.} The minimal witness is a single
read--write antidependency from $i$ to $j$, serializable in the order $i;j$,
and none of Adya's proscribed phenomena~\cite{Adya1999}; TLC exhibits such a
history (claim \texttt{A1OnlyOnNonSerializable}, violated as expected). \A1
is therefore not an isolation anomaly in the classical sense. Its
\emph{absence} is nevertheless a classical correctness class, and we state
which one at the level at which the theorem holds.

\begin{definition}[Commit-order read consistency]
\label{def:corc}
A history is \emph{commit-order read consistent} when every operation read, for
every cell in its read set, exactly the value it would have read had it
executed atomically at its own commit point: for all $i$ and $c\in h[i].\rs$,
$h[i].\rv[c]$ is the value of the latest write to $c$ with
$\wt<h[i].\wt$, or $\mathrm{NULL}$ if there is none.
\end{definition}

Writes are installed in commit order in any case, so \cref{def:corc} is
value-level equivalence, on the reads, to the serial execution in commit
order. It is the value-level analogue of Weikum and Vossen's
commit-order-preserving serializability~\cite{WeikumVossen2001} and not that
condition itself, which is conflict-based; under value-unique writes the two
coincide on reads (view equivalence~\cite{Papadimitriou1979}).

\begin{theorem}[\A1-freedom implies commit-order read consistency]
\label{thm:corc}
Let $h$ satisfy the store axiom, read-before-commit and per-agent non-overlap.
If $h$ has no \A1 witness, then $h$ is commit-order read consistent.
\end{theorem}

The proof (\texttt{CommitOrder.tla}; TLAPS, \TlapsCommitOrder{} obligations,
no omitted step) is Kung and Robinson's backward-validation argument run
against our definitions: if some operation's read of $c$ disagrees with the
latest write to $c$ before its commit, that write lies inside the operation's
window, was made by another agent (one agent's operations do not overlap), and
wrote a different value---an \A1 witness. Fifteen TLC verdicts
(\texttt{check\_commitorder.sh}) fix what the statement says, and all fifteen
reproduce: over all \CommitOrderHistories{} histories of length at most three
the implication holds; it is not vacuous; it fails when either droppable
hypothesis is removed; the relational definitions equal
\texttt{Memory.tla}'s and a literal serial replay; the three hypotheses are
invariants of the operational model over \CommitOrderMemoryStates{} reachable
states; and the converse fails in exactly one shape---a value overwritten and
then restored inside the window fires \A1 although the history is
commit-order read consistent. The predicates nest strictly:
\[
  \neg\text{CORC} \;\Longrightarrow\; A_1 \;\Longrightarrow\; \text{validation fails}.
\]
\emph{Reach.} The theorem is a conditional guarantee over faithful histories
that satisfy its three hypotheses; it says nothing about serializability of
an application, about effects, about reads the instrumentation did not
declare, or about a runtime that violates per-agent non-overlap. One deployed
construct does: LangGraph's \texttt{Send} fan-out runs several instances
under one agent name and is outside the model (\cref{sec:model}). On the
committed data both droppable hypotheses hold---\HypSameAgentOverlaps{}
same-agent overlapping windows and \HypReadBeforeCommitViolations{} reads
not before their commit in \HypRecords{} records over \HypTraces{} traces
(\texttt{hypothesis\_gates.py}, a gate that fails on the first
violation)---and a runtime can enforce the first by refusing a second
operation of an agent whose window is open. The converse shape is a false positive of $L_1$, priced in
discarded generations; on the live records it occurs in \HypRestored{} of
\HypAOneWitnesses{} \A1 witness pairs.
Two corollaries are mechanized with the theorem. The classical lost update
(P4 of Berenson et al.~\cite{BerensonEtAl1995}) is an \A1 witness in which
the stale reader also writes the cell it read, so $L_1$ excludes it. A
\emph{blind} overwrite---a write by an operation that never read the cell---is
not an anomaly in isolation theory and is outside the catalog; it matters in
\cref{sec:prevalence}, because the one in-the-wild bug we reproduce is of that
kind, and in \cref{sec:live-a1}, because a revision that lands \emph{after}
the reader's commit is of that kind too.

\subsection{\A2: phantom tool}
\label{sec:a2}
\begin{definition}[Phantom tool]
\label{def:a2}
$h$ exhibits a \emph{phantom tool} when for some $i$,
\[
  \begin{gathered}
  h[i].\pt\neq\mathrm{NULL},\\
  h[i].\pt\in h[i].\rreg,\\
  h[i].\pt\notin h[i].\wreg.
  \end{gathered}
\]
\end{definition}
An agent selects a tool from the registry it read; the registry mutates
during generation; the commit emits a call to a capability that no longer
exists. TLC's witness has four states. \emph{Tool validity}, canonically, is
membership of the (identifier, signature) pair: a tool is the same tool while
both are unchanged. \Cref{def:a2} is stated over identifiers, the \TLA{}
registry being a set of them; the Verus restatement of \cref{sec:l4} maps
identifiers to signatures and fires additionally when a tool is re-signed
under its old name. The two coincide under the encoding the executable
dispatcher uses, in which re-signing a tool is removal of the old pair and
insertion of a fresh one, so that a re-signed tool is a removed tool to
\cref{def:a2}; the scorer and the $L_4$ dispatcher decide the pair form, and
\cref{tab:assurance} records the correspondence as an encoding, not a
proof. A tool the agent names that was never in $\rreg$ (a hallucinated
call) is not \A2 and not in the catalog.

\subsection{\A3: causal cascade}
\label{sec:a3}
When an operation is retracted in a compensating (saga-style) runtime, any
operation that read a value it produced is grounded on a basis no surviving
commit supports. The classical remedy, cascading abort, flags the dependent
aborted. In the agent regime that remedy fails in a specific way: if the
dependent's tool effects have already been issued, the flag changes nothing
that reached the world. The cataloged predicate is therefore stated over
externalization.

\begin{definition}[Causal cascade]
\label{def:a3}
Let each record carry $\fld{aborted}$ and $\fld{externalized}$ flags and a
predecessor set $\fld{preds}$ (the operations whose committed writes it
observed, transitively closed). $h$ exhibits a \emph{causal cascade} when for
some $j,p$:
\[
  h[j].\fld{externalized}\;\wedge\;p\in h[j].\fld{preds}\;\wedge\;h[p].\fld{aborted}.
\]
\end{definition}
The surviving-dependent form, $\neg h[j].\fld{aborted}\wedge p\in
h[j].\fld{preds}\wedge h[p].\fld{aborted}$, is kept beside it as
\textsc{CascadeUnpropagated}; a cascading abort falsifies that form by
flagging $j$ after $j$'s effects are out, and \cref{sec:l2-eval} measures the
separation on runnable arms.

\smallskip\noindent\emph{Prevention is output commit, and output commit has a
residue.} Under \cref{def:a3} a runtime must hold an operation's effects until
every operation in its closure has released its own, and must refuse to
retract an operation whose effects are out. The first clause is prevention.
The second keeps the predicate false by refusing to set the flag it reads, and
we do not count that as prevention: a refused retraction is a basis someone
wished to withdraw after release, the ordinary rule of recoverability that a
committed transaction cannot abort, and every measurement in
\cref{sec:l2-eval} counts its residue beside the zero.

\smallskip\noindent\emph{What the detector decides.} A trace without
provenance fields cannot decide \cref{def:a3}. What it can decide is an
\emph{unsupported read}---a committed read of a value that no surviving
committed write produced at or before the read. The verified
\texttt{detect\_a3} decides exactly that screen, and its relation to
\cref{def:a3} is five mechanized facts (supplement \S\,G): it is not sound
in general (a value collision hides the aborted writer), it is sound under
value-unique writes, the hypothesis black-box isolation checkers also
need~\cite{Aphyr2020Elle,TanEtAl2020Cobra}, it holds along chains, it is not
complete, and it is blind to relabeled effects. We therefore claim
\cref{def:a3} only where the hypothesis holds, and check it: value-unique
writes hold by construction in the synthetic cascades (distinct value ranges
per writer) and in all \HypLiveLtwoSessions{} live replay sessions
(\HypLiveLtwoCollisions{} collisions), the two data sets the detector scores;
they fail in \HypCollisionCells{} cells of the pilot and occurrence-study
traces, where only \A1 is scored. A recorder that stamps each written value with its writer's operation
identifier makes the hypothesis a system invariant.

\subsection{\A6: tool-effect reordering}
\label{sec:a6}
Each record additionally carries an issuance order $\fld{io}$ and a commit
order $\fld{co}$ over its external effects; $\ws$ and $\wv$ derive from
$\fld{co}$.
\begin{definition}[Tool-effect reordering]
\label{def:a6}
$h$ exhibits \emph{tool-effect reordering} when for some $i$,
\[
  \begin{gathered}
  |h[i].\fld{io}|\ge 2,\quad h[i].\fld{co}\neq h[i].\fld{io},\\
  h[i].\fld{co}\ \text{is a permutation of}\ h[i].\fld{io}.
  \end{gathered}
\]
\end{definition}
The minimal witness (three states, one agent, two cells) is benign over store
cells, where writes to distinct cells commute; the predicate names an anomaly
only under the reading in which each element of $\fld{io}$ is an external
effect whose order an observer can see. It over-reports on commuting effects
and, because it requires a permutation, is silent on dropped and duplicated
effects; \cref{sec:limitations} exhibits the duplicated case in a deployed
framework.


\subsection{Snapshot insufficiency}
\label{sec:snapshot}
$L_1$ is not a snapshot guarantee: ``reads reflect memory at read time''
is the store axiom of \cref{thm:corc}, holds on every reachable log, and
TLC exhibits a history satisfying it on which \A1 fires
(\texttt{MC\_A1\_struct}; supplement \S\,B). The live instance is a runtime
that pins a consistent snapshot per node---LangGraph's superstep---and has
not thereby obtained $L_1$ (\cref{sec:live-a1}).

\subsection{What the catalog leaves: the residue predicate}
\label{sec:residue}
Every predicate above is decided at or before an operation's commit, and a
discipline that excludes all four still leaves a harm the live study measures
in a third of sessions (\cref{sec:live-a1}): the external effect a commit
released is irreversible, while the state it was derived from is ordinary
revisable state that any later operation may overwrite. This is not a
concurrency anomaly. It fires in correct sequential executions---book
against the date known at commit, then be contradicted by a legitimate
revision---and classical isolation has no name for it because a committed
transaction's inputs are final by definition; the nearest notions are view
maintenance~\cite{GuptaMumick1995ViewMaint} with the view outside the system,
and the saga's compensation~\cite{GarciaMolinaSalem1987Sagas}. We define it because an evaluation scored on outcomes needs a predicate
for the outcome's other half, and keep it outside the catalog and the
lattice.

Each operation record carries a \emph{basis}, the (cell, value) pairs its
external effect was derived from, and an $\fld{obligation}$ flag. The basis
is recovered by \emph{argument provenance}: a read value that appears
verbatim among the effect's arguments is part of its basis, and nothing is
declared by the agent. This is why-provenance~\cite{BunemanEtAl2001Why}
restricted to copying, with that tradition's known limit: a value that
reaches the effect through a computation leaves no basis, so the extractor
under-reaches and never over-reports. In the model the basis is a
nondeterministic subset of the read values, so every extractor is covered.

\begin{definition}[Residue predicate]
\label{def:residue}
$h$ exhibits a \emph{superseded-basis effect} (\Rp{}; \texttt{NoA7} in the
artifact) when for some $i$, $h[i].\fld{externalized}$,
$\neg h[i].\fld{obligation}$, and for some $(c,v)\in h[i].\fld{basis}$ the
\emph{latest} write to $c$ committed after $h[i].\wt$ wrote a value other
than $v$:
\[
  \begin{gathered}
  \exists j:\ c\in h[j].\ws \;\wedge\; h[j].\wt > h[i].\wt\\
  \qquad\wedge\; h[j].\wv[c]\neq v\\
  \wedge\ \forall k:\ (c\in h[k].\ws \wedge h[k].\wt > h[i].\wt)\\
  \qquad\qquad\Rightarrow h[k].\wt \le h[j].\wt.
  \end{gathered}
\]
\end{definition}

\A1 and \Rp{} partition ``an effect is out on a basis the state no longer
holds'' by whether the superseding write committed inside the operation's
window or after its commit; a value overwritten and then restored fires
neither, and \Rp{} carries no agent conjunct because the obligation is owed
whoever supersedes. On the live data the partition is exact
(\cref{tab:live}): $A_1 + R = 100$ in every unguarded cell, and under $L_1$
the \A1 column is empty and the residue is entirely \Rp{}. The two sides of
that identity are computed independently---the outcome column compares the
booking with the final store state and never reads a basis; \Rp{} reads the
trace's write order and never the final state---so the identity is a check
on the extractor, not a tautology of it; on this workload the extractor
recovers the basis of \HypBasisRecovered/\HypBookings{} bookings because the
booking tool takes the date as a structured argument, and its blind spot is
stated in \cref{sec:limitations}. TLC's witness has five states and is not
an \A1 witness: an operation reads $c_1=\mathrm{NULL}$ and commits an effect
with basis $\{(c_1,\mathrm{NULL})\}$, released at commit; a later operation
commits $c_1=v_1$.

\begin{proposition}[What remains after release]
\label{prop:owed}
Let $D$ be any discipline under which an effect $e$ with basis $B$ is
released at a finite time $T$, and let the store accept, at any time after
$T$, a write to a cell of $B$ by an operation that does not depend on $e$
(the open-world hypothesis). (i)~There is an execution of $D$ in which a
write committed after $T$ supersedes $B$ while $e$ is out. (ii)~If
obligations are never cleared, $\neg R$ holds in every reachable state of
$D$ if and only if, in every execution, $e$'s obligation is set in the same
step as, or before, the step that completes the other conjuncts of
\cref{def:residue}---the superseding commit for a released effect, the
release for an effect superseded while held.
\end{proposition}
Both halves are immediate---(i) because the write is admitted after a
release that cannot be recalled, (ii) because \cref{def:residue} is
``superseded and not owed'' and a monotone flag set by that step stays
set---and we record them as a proposition, not a theorem, because their
content is to fix what a discipline can promise about \Rp{} under the
open-world hypothesis: not a current basis, but an obligation registered
atomically with the step that completes the witness, which the ledger of
\cref{sec:ledger} does at both steps. Dropping the hypothesis is possible
and is not free: freezing the cells of $B$ until the effect is acknowledged
is a lock held past commit, with the pessimistic column of \cref{tab:cost}
as its price floor; output commit (\cref{sec:a3}) holds effects only until
their \emph{causal} basis is externalized, a finite condition. What a
registered obligation buys in the world---a rebooking, a cancellation, a
human notification---is a property of the compensator and is measured, not
proved (\cref{sec:live-a1}).
